\documentclass[11pt,a4paper]{article}

\usepackage[T1]{fontenc}
\usepackage[utf8]{inputenc}
\usepackage[margin=2.4cm]{geometry}
\usepackage{lmodern}
\usepackage{amsmath,amssymb}
\usepackage{listings}
\usepackage{xcolor}
\usepackage{enumitem}
\usepackage{fancyhdr}
\usepackage[hidelinks]{hyperref}
\usepackage{microtype}
% Command-line flags must render as two hyphens, not an en-dash.
\DisableLigatures[-]{encoding=T1,family=tt*}

\definecolor{codebg}{gray}{0.96}
\definecolor{codekw}{RGB}{20,80,160}
\definecolor{codecm}{RGB}{100,110,100}
\definecolor{codestr}{RGB}{150,60,20}

\lstdefinestyle{src}{
  backgroundcolor=\color{codebg},
  basicstyle=\ttfamily\footnotesize,
  keywordstyle=\color{codekw}\bfseries,
  commentstyle=\color{codecm}\itshape,
  stringstyle=\color{codestr},
  showstringspaces=false,
  columns=fullflexible,
  keepspaces=true,
  breaklines=true,
  frame=leftline,
  framerule=1.2pt,
  rulecolor=\color{codekw!40},
  xleftmargin=10pt,
  framexleftmargin=8pt,
  aboveskip=8pt,
  belowskip=8pt,
  literate={°}{{\textdegree}}1 {→}{{$\rightarrow$}}1,
}
\lstdefinestyle{sh}{
  style=src,
  language=bash,
  morekeywords={esbmc,ltl2ba},
}
\lstset{style=src,language=C}

\newcommand{\esbmc}{ESBMC}
\newcommand{\flag}[1]{\texttt{\small #1}}
\newcommand{\file}[1]{\texttt{\small #1}}

\newlist{parts}{enumerate}{1}
\setlist[parts]{label=(\alph*), leftmargin=2em, topsep=4pt, itemsep=3pt}

\newcounter{exc}
\renewcommand{\theexc}{\arabic{exc}}
\makeatletter\newcommand{\excref}[1]{\def\@currentlabel{\theexc}\label{#1}}\makeatother
\newenvironment{exercise}[2][]
  {\stepcounter{exc}\ifx\relax#1\relax\else\excref{#1}\fi%
   \medskip\noindent\textbf{Exercise \theexc\ (#2)}\par\nopagebreak\smallskip}
  {\par\medskip}

\pagestyle{fancy}
\fancyhf{}
\lhead{\small Specification and Verification of Embedded and Cyber-Physical Systems}
\rhead{\small \thepage}
\renewcommand{\headrulewidth}{0.4pt}

\title{\vspace{-1.5cm}Specification and Verification of\\Embedded and Cyber-Physical Systems\\[4pt]
       \large Exercise List --- \esbmc{} 8.5}
\author{}
\date{}

\emergencystretch=3em

\begin{document}
\maketitle
\vspace{-1.6cm}

\section*{Before you start}

Every command in this list was run against \esbmc{} \textbf{8.5.0}. Check your
build first:

\begin{lstlisting}[style=sh]
esbmc --version      # expect 8.5.0 or newer
\end{lstlisting}

\noindent
Reference [3] is the tool paper; [1] describes the SMT encoding underneath it.

Three things changed since the version used in earlier editions of this
tutorial, and they will show up in your terminal:

\begin{itemize}[leftmargin=1.4em,itemsep=2pt]
  \item The default solver is now \textbf{Bitwuzla}. Use \flag{--z3},
        \flag{--cvc5} or \flag{--boolector} to cross-check a result.
  \item Verification \textbf{stops at the first violated property}; the rest
        are listed as \texttt{NOT CHECKED}. Add \flag{--multi-property} for a
        verdict on every property in a single run.
  \item Loop invariants in C are written \flag{\_\_ESBMC\_loop\_invariant(...)}
        and enabled with \flag{--loop-invariant} or
        \flag{--loop-invariant-check}. The Python frontend spells the same
        annotation \flag{\_\_loop\_invariant(...)}.
\end{itemize}

Unless stated otherwise, the files referred to below live in \file{examples/}
in the tutorial repository; \file{bmc-example.c} and \file{bmc-example.py} are
in \file{exercises/}, next to this handout. \file{examples/README.md} records
the command and the expected verdict for each example, and
\file{examples/run-all.sh} re-checks them all.

\medskip
\noindent
\textbf{How to answer.} Where an exercise asks you to run \esbmc{}, report the
command, the verdict, and --- when the verdict is \texttt{FAILED} --- the
counterexample values that make the property fail. A verdict on its own is not
an answer; the interesting part is always \emph{why}.

\medskip
\noindent
\textbf{Runtimes.} Most runs finish in under a second. The concurrent ones in
Sections 4 and 6 take from a few seconds to a few minutes and grow sharply with
\flag{--context-bound} and \flag{--unwind}; raise them one step at a time.

\section{Bounded model checking}

\begin{exercise}{Bounded Model Checking}
Describe the steps a bounded model checker takes, from reading the program to
producing an SMT formula, using the program below (see [1] for the encoding
\esbmc{} actually uses). The \texttt{\_\_ESBMC\_assume}
directive constrains non-deterministic variables; \texttt{assert} states the
property to check.

\begin{lstlisting}
#include <assert.h>
#define a 2
int main(int argc, char **argv) {
  long long int i = 1, sn = 0;
  unsigned int n = 5;
  __ESBMC_assume(n >= 1);
  while (i <= n) {
    sn = sn + a;
    i++;
  }
  assert(sn == n * a);
}
\end{lstlisting}

\begin{parts}
  \item Unroll the loop by hand for $n = 5$ and write the resulting
        single-assignment (SSA) form.
  \item Split the SSA form into the constraint $C$ and the property $P$, and
        give the formula the solver is asked to decide. Which answer ---
        \emph{satisfiable} or \emph{unsatisfiable} --- corresponds to
        \texttt{VERIFICATION FAILED}?
  \item Run \texttt{esbmc bmc-example.c}. Compare the number of assignments
        reported by \emph{Symex completed} with the number that survives
        \emph{Slicing}. What has the slicer removed, and why is it sound to
        remove it?
  \item Replace \texttt{unsigned int n = 5;} with
        \texttt{unsigned int n = nondet\_uint();}. The verdict changes. Explain
        what \flag{--unwind} now controls, and what an
        \emph{unwinding assertion} is telling you when it fails.
\end{parts}
\end{exercise}

\begin{exercise}{Completeness of BMC}
BMC is a bug-finding technique first and a proof technique second.

\begin{parts}
  \item State the condition under which \texttt{VERIFICATION SUCCESSFUL} from a
        bounded run is a \emph{proof} that the property holds for all
        executions.
  \item Run \file{ex1.cpp} at \flag{--unwind 5} twice: once with
        \flag{--no-unwinding-assertions} and once without. Both report
        \texttt{VERIFICATION FAILED}, but they fail on different properties ---
        name them. Say which of the two runs you would be willing to put in a
        safety argument, and note that ``\texttt{FAILED}'' on its own told you
        almost nothing here.
  \item The Python program \file{collatz} below terminates for
        \texttt{collatz(6)} but not for \texttt{collatz(-1)}. Explain what a
        bounded model checker can and cannot tell you about it.

\begin{lstlisting}[language=Python]
def collatz(n: int) -> int:
    while n != 1:
        if n % 2 == 0:
            n = n // 2
        else:
            n = 3 * n + 1
    return 1
\end{lstlisting}
\end{parts}
\end{exercise}

\section{Loops: invariants, k-induction and abstract domains}

\begin{exercise}{Strengthening a loop invariant}
\file{examples/incomplete.c} carries two loop invariants that are deliberately
too weak:

\begin{lstlisting}
  __ESBMC_loop_invariant(i >= 0);
  __ESBMC_loop_invariant(j >= i);
  while (i <= n) { i = (i + 1); j = (j + i); }
  assert((i + (j + k)) > (2 * n));
\end{lstlisting}

\begin{parts}
  \item Run \texttt{esbmc incomplete.c --loop-invariant --no-standard-checks}.
        Which of the three obligations --- \emph{base case}, \emph{inductive
        step}, \emph{postcondition} --- fails, and what does that tell you about
        where the invariant is too weak?
  \item Strengthen the annotations until the run reports
        \texttt{VERIFICATION SUCCESSFUL}. Three things are needed, and the
        second and third are the interesting ones:
        \begin{enumerate}[label=\roman*.,leftmargin=2.2em,topsep=2pt]
          \item an upper bound on $i$ that pins its value at loop exit;
          \item an invariant that fixes $j$ \emph{exactly} rather than bounding
                it. A range such as $j \ge i$ is inductive but too weak: the
                inductive step havocs the loop variables, so it admits states
                the loop can never actually reach. Read the counterexample
                \esbmc{} prints and find the havocked state it used;
          \item bounds on the inputs \texttt{n} and \texttt{k}. Without them the
                run still fails --- read that counterexample too and say which
                expression overflows.
        \end{enumerate}
        State the final annotations and argue informally that the invariant is
        inductive.
  \item Repeat on the Python version, \file{examples/incomplete.py}, with
        \flag{--loop-invariant-check}. The annotation is spelled
        \flag{\_\_loop\_invariant} there. All three of the things above are
        needed again, including the input bounds --- say what that tells you
        about the integer model \esbmc{}'s Python frontend uses, and how it
        differs from CPython's.
  \item \file{examples/crosshair.py} is a solved instance of the same shape.
        Explain why the equalities $y = i$ and $x = 2i$ are needed and why the
        bounds $0 \le i \le n$ alone would not be enough.
\end{parts}
\end{exercise}

\begin{exercise}{k-induction and interval analysis}
\begin{parts}
  \item Run \texttt{esbmc kinduction-interval.c --k-induction}. The verdict is
        \texttt{VERIFICATION UNKNOWN}. Read the output and say which of the
        three k-induction steps --- base case, forward condition, inductive
        step --- did not go through.
  \item Add \flag{--interval-analysis} and run it again. It now succeeds. What
        fact about \texttt{status} does the interval domain supply that
        k-induction could not derive on its own?
  \item Do the same for \file{kinduction-interval2.c}, where the loop counter
        starts non-deterministically.
  \item \file{kinduction.c} succeeds with plain \flag{--k-induction}, and the
        output reports the $k$ at which it closed. Explain what
        ``\emph{Solution found by the inductive step ($k = 2$)}'' means
        operationally, and relate the three steps \esbmc{} runs --- base case,
        forward condition, inductive step --- to the algorithm in [2].
\end{parts}
\end{exercise}

\begin{exercise}{Three ways to handle the same loop}
The three Python files \file{loop-invariant.py}, \file{loop-invariant2.py} and
\file{loop-invariant3.py} attack one loop in three ways: full unwinding,
truncation to one iteration, and replacement by an invariant.

\begin{parts}
  \item Run all three (\file{loop-invariant.py} needs \flag{--unwind 101}) and
        record the verdicts.
  \item \file{loop-invariant2.py} reports \texttt{VERIFICATION FAILED} even
        though the original program is correct. Name the phenomenon this
        illustrates and say which way it errs --- is the tool reporting a bug
        that cannot happen, or missing one that can? Then say which way a
        bounded run errs when \flag{--no-unwinding-assertions} is passed. The
        two answers are opposites, and only one of them is unsoundness.
  \item \file{loop-invariant3.py} succeeds without unwinding the loop at all.
        Which of its \texttt{\_\_ESBMC\_assume} lines encode the invariant, and
        which encodes the negated loop condition?
  \item Which of the three approaches scales to a loop bound of $10^9$, and at
        what cost to the engineer?
\end{parts}
\end{exercise}

\section{Memory safety and arithmetic}

\begin{exercise}{Memory safety}
\begin{parts}
  \item \file{examples/ex1.cpp} has \emph{two} defects, both on line 7. Find
        them by inspection, then run
        \texttt{esbmc ex1.cpp --unwind 5 --no-unwinding-assertions}. Only one is
        reported. Add \flag{--multi-property}, list all four properties \esbmc{}
        raises on that line, and say which two hold and why.
  \item \esbmc{} distinguishes \emph{array bounds violated},
        \emph{NULL pointer}, \emph{invalidated dynamic object} and
        \emph{incorrect alignment}. Give a one-line C fragment that triggers
        each, and verify each with \esbmc{}.
  \item Fix \file{ex1.cpp} and show that the fixed version verifies for an
        unbounded \texttt{n}. Plain \flag{--k-induction} returns
        \texttt{VERIFICATION UNKNOWN} here, with or without
        \flag{--interval-analysis}; a loop invariant relating the counter to
        \texttt{n} does discharge it. State the invariant and the flags you
        used.
\end{parts}
\end{exercise}

\begin{exercise}{Arithmetic and the number model}
\begin{parts}
  \item Run \texttt{esbmc add-overflow.py}. It succeeds. Now add
        \flag{--overflow-check}. Explain, in terms of verification conditions,
        why the first run could not have found the bug.
  \item Run \texttt{esbmc neural-net.c} and then
        \texttt{esbmc neural-net.c --fixedbv}. The same assertion passes under
        one number model and fails under the other. Explain which model each
        flag selects and which of the two verdicts describes the real hardware.
  \item \file{neural-net.py} computes the same network and succeeds. What does
        that tell you about comparing verification results across language
        frontends? See [4] for how the Python frontend is built.
\end{parts}
\end{exercise}

\section{Concurrency}

\begin{exercise}{Interleavings and data races}
\begin{parts}
  \item \file{examples/interleavings.c} runs \texttt{x++} twice against
        \texttt{x--} twice with no synchronisation. Enumerate the values
        \texttt{x} can hold once both threads have run to completion, and
        justify each with an interleaving of the load/add/store steps. Note
        first that \texttt{main} never joins the threads, so the program as
        written has no such point --- say what that means for the question, and
        add the two \texttt{pthread\_join} calls before answering it.
  \item Run \texttt{esbmc interleavings.c --data-races-check}. Note that
        \esbmc{} reports the race without any assertion in the program. What
        is the difference between a \emph{data race} and an
        \emph{assertion violation}?
  \item Read the \emph{Schedule reduction} line at the end of the output. Say
        what \texttt{peak\_threads}, \texttt{schedules\_explored} and the four
        \texttt{pruned\_by\_*} counters report. On this example every pruning
        counter is \emph{zero} --- explain why partial-order reduction finds
        nothing to prune here, and construct a variant of the program on which
        \texttt{pruned\_by\_mpor} is non-zero.
\end{parts}
\end{exercise}

\begin{exercise}{Atomicity violation in a shared bank account}
\file{examples/bank-account-bug.c} checks the balance and then withdraws, with
the two steps outside any lock. \file{examples/bank-account.c} is the same code
with a mutex around both.

\begin{parts}
  \item Run both with \texttt{--unwind 2 --no-unwinding-assertions}. Report the
        two verdicts and, for the failing one, the counterexample values of
        \texttt{balance}.
  \item The bug is not a data race in the strict sense --- each individual
        access could be atomic and the bug would remain. Explain why, and name
        the class of defect this is.
  \item Suppose the \texttt{pthread\_mutex\_unlock} in \file{bank-account.c}
        were moved to just after the \texttt{if (balance >= amount)} test.
        Explain what goes wrong, and give the interleaving of the two threads
        that overdraws the account. Then explain why this variant is far more
        expensive for \esbmc{} to check than either of the two originals ---
        it does not finish within a few minutes at these bounds --- in terms of
        how many context switches the split critical section admits.
\end{parts}
\end{exercise}

\begin{exercise}{Mutual exclusion}
Consider Peterson's algorithm as implemented in
\file{examples/mutual-exclusion.c}.

\begin{lstlisting}
void *t1(void *arg) {                  void *t2(void *arg) {
  flag[0] = 1;                           flag[1] = 1;
  turn = 1;                              turn = 0;
  while (flag[1]==1 && turn==1) {};      while (flag[0]==1 && turn==0) {};
  c1 = 1;                                c2 = 1;
  CT0: if (i==1) x=1;                    CT1: if (i==2) x=2;
  c1 = 0;                                c2 = 0;
  flag[0] = 0;                           flag[1] = 0;
  return NULL;                           return NULL;
}                                      }
\end{lstlisting}

\begin{parts}
  \item Specify the following three properties in LTL, using $c_i$ for
        ``process $i$ is in its critical section'' and $r_i$ for ``process $i$
        has requested entry'' ($i \in \{1,2\}$; the propositions are named
        apart from the functions \texttt{t1} and \texttt{t2} on purpose).
        \begin{enumerate}[label=\roman*.,leftmargin=2.2em,topsep=2pt]
          \item At most one process is in the critical region at any time.
          \item Whenever a process tries to enter its critical region, it will
                eventually succeed.
          \item A process can eventually ask to enter its critical region.
        \end{enumerate}
  \item Which of the three are \emph{safety} properties and which are
        \emph{liveness} properties (see [6])? Which of them can a
        bounded run refute, and which can it only fail to refute?
  \item Verify property (i):
\begin{lstlisting}[style=sh]
esbmc mutual-exclusion.c --unwind 3 --context-bound 3 \
      --no-unwinding-assertions
\end{lstlisting}
  \item Show that the critical section is reachable at all, by checking the
        label \texttt{CT0} with \flag{--error-label CT0}. Why is a bare
        \flag{--error-label} verdict \emph{not} on its own evidence that a label
        is unreachable?
  \item Delete \texttt{turn = 1;} from \texttt{t1} and re-verify. Report the
        interleaving \esbmc{} produces.
\end{parts}
\end{exercise}

\begin{exercise}{Locking a counter}
\file{examples/pthread.c} has two threads increment a shared \texttt{n} ten
times each under a mutex, and asserts \texttt{n == 20}.

\begin{parts}
  \item Verify it with \texttt{--unwind 11 --context-bound 2}.
  \item Remove the \texttt{pthread\_mutex\_lock}/\texttt{unlock} pair and
        re-verify at the same bounds (\flag{--unwind 11}
        \flag{--context-bound 2}). What is the smallest final value of
        \texttt{n} that \esbmc{} can justify with an interleaving? The answer
        depends on the context bound --- say what it would become without one.
  \item Raise \flag{--context-bound} and observe the runtime. State precisely
        what a run with \flag{--context-bound}~$k$ does \emph{not} prove.
\end{parts}
\end{exercise}

\section{Temporal logic}

\begin{exercise}{Specifying and checking an LTL property}
\esbmc{} checks LTL by compiling the \emph{positive} formula into a Büchi
monitor with the ESBMC fork of \texttt{ltl2ba} and running it as an extra
thread. The tool performs the negation internally.

\begin{lstlisting}[style=sh]
ltl2ba -O c -H '"vars.h"' -f 'G({pressed} -> F {charge > min})' \
       > ltl_example.ba-2.c
esbmc ltl_example.c --ltl ltl_example.ba-2.c -DLTL_PREFIX_BOUND=10
\end{lstlisting}

\begin{parts}
  \item Run the two commands above on \file{examples/ltl\_example.c} and report
        the \emph{Final lowest outcome}. \esbmc{} distinguishes
        \texttt{LTL\_BAD}, \texttt{LTL\_FAILING} and \texttt{LTL\_SUCCEEDING} ---
        explain what each means for a \emph{finite prefix} of an infinite trace,
        using the good-prefix/bad-prefix vocabulary of [6].
  \item Explain the role of \texttt{LTL\_PREFIX\_BOUND}. Reduce it to 1 and
        explain the resulting verdict.
  \item Reduce the loop in \file{ltl\_example.c} to a single iteration. The
        verdict becomes \texttt{VERIFICATION UNKNOWN}. Why does the monitor need
        more than one step to reach a decision?
  \item Write and check the LTL formula
        $G(\text{pressed} \rightarrow X\,(\text{charge} > \text{min}))$. Is it
        stronger or weaker than the formula above? Confirm your answer with
        \esbmc{}.
\end{parts}
\end{exercise}

\begin{exercise}[ex:plant]{Specifying a chemical-process controller}
A typical embedded system consists of a human--machine interface (keyboard and
display), a processing unit, and an instrumentation interface (sensors,
network, actuators) connected to a physical plant. The system considered here
controls a chemical process and can be operated remotely over TCP/IP. It has
two ADCs --- one reading thermocouples, one reading a pressure transducer ---
and a DAC that sets the valve/pump position, plus a switch for the heater. Its
objective is to keep temperature and pressure within defined limits.

\begin{parts}
  \item Specify each of the following in LTL:
        \begin{enumerate}[label=\alph*.,leftmargin=2.2em,topsep=2pt]
          \item A T process reads a temperature sensor, turns the heating system
                on if the temperature is below 20\,\textdegree C, and turns it
                off if the temperature is above 300\,\textdegree C.
          \item A P process regulates pressure, opening the pump/valve when
                pressure is above 500\,bar and closing it when pressure is below
                100\,bar.
          \item The measured values are shown on an LCD when the T and P
                processes transfer data to an S process.
        \end{enumerate}
  \item Which of your three formulas are safety properties? For each safety
        property, give the corresponding C assertion.
  \item What are the security objectives of this system?
  \item What are the sources of security problems that arise when it is operated
        remotely? Consider the TCP interface specifically.
\end{parts}
\end{exercise}

\section{Verifying LLM-generated code}

\begin{exercise}{Reviewing an LLM implementation with a model checker}
The directory \file{ai-generated/} contains three LLM-written implementations of
the controller from Exercise~\ref{ex:plant}: \file{embedded\_system\_llm.c} (no
synchronisation), \file{embedded\_system\_unsafe1.c} (one global mutex) and
\file{embedded\_system\_unsafe2.c} (one mutex per shared variable).

\begin{lstlisting}[style=sh]
esbmc embedded_system_llm.c     --data-races-check --context-bound 2 \
      --unwind 2 --no-unwinding-assertions
esbmc embedded_system_unsafe2.c --data-races-check --context-bound 2 \
      --unwind 2 --no-unwinding-assertions
esbmc embedded_system_unsafe1.c --data-races-check --context-bound 1 \
      --unwind 2 --no-unwinding-assertions
\end{lstlisting}

The first two take one to two minutes; the third takes a few seconds. Note the
lower context bound on the third --- at \flag{--context-bound 2} it does not
finish in several minutes.

\begin{parts}
  \item Run all three. Report the verdict for each and, where one fails, the
        first property violated.
  \item The third file reports \texttt{VERIFICATION SUCCESSFUL}, yet it takes
        the same non-recursive mutex twice on the path
        \texttt{temperature\_process} $\rightarrow$ \texttt{set\_heater}, which
        self-deadlocks under POSIX. Adding \flag{--deadlock-check} does not
        change the verdict either. Explain why the bounded run cannot see this
        defect --- look at what the worker threads do, and at where \esbmc{}
        places its deadlock property --- and propose a change to the harness
        that would expose it.
  \item In \file{embedded\_system\_unsafe2.c} the assertion relating
        \texttt{temperature} to \texttt{heater\_status} was moved into the
        display process, and each variable has its own lock. Explain why
        per-variable locking does not preserve the relation between them.
  \item All three files were produced by a language model from the same natural
        language specification, and all three look plausible. Write two or three
        sentences on what this implies for reviewing generated code, and on what
        a model checker adds that a human reviewer does not. Reference [5]
        describes one way to close the loop: feed the counterexample back to the
        model and re-verify the repair.
  \item Repair one of the three so that it verifies, and state the property your
        repaired version establishes and the bounds under which it establishes
        it.
\end{parts}
\end{exercise}

\begin{exercise}{Writing a specification an LLM cannot fake}
\file{examples/llm\_thermostat\_example.c} is a short LLM-written guard that
\esbmc{} reports as \texttt{VERIFICATION SUCCESSFUL}.

\begin{parts}
  \item Read the program and explain why the successful verdict is close to
        meaningless.
  \item Rewrite it so that the temperature is non-deterministic and the
        assertion states a property of the \emph{controller} rather than of a
        constant. Verify your version.
  \item Formulate one property of a thermostat that BMC cannot establish at all,
        and say what technique you would use instead.
\end{parts}
\end{exercise}

\section*{References}

\sloppy
\begin{enumerate}[leftmargin=2.4em,itemsep=3pt,label={[\arabic*]}]
  \item L. C. Cordeiro, B. Fischer, J. Marques-Silva. ``SMT-Based Bounded Model
        Checking for Embedded ANSI-C Software''. \emph{IEEE Transactions on
        Software Engineering} 38(4):957--974, 2012.
        \href{https://doi.org/10.1109/TSE.2011.59}{doi:10.1109/TSE.2011.59}
  \item M. Y. R. Gadelha, H. I. Ismail, L. C. Cordeiro. ``Handling loops in
        bounded model checking of C programs via k-induction''.
        \emph{International Journal on Software Tools for Technology Transfer}
        19(1):97--114, 2017.
        \href{https://doi.org/10.1007/s10009-015-0407-9}{doi:10.1007/s10009-015-0407-9}
  \item M. Y. R. Gadelha, F. R. Monteiro, J. Morse, L. C. Cordeiro, B. Fischer,
        D. A. Nicole. ``ESBMC 5.0: an industrial-strength C model checker''.
        \emph{ASE 2018}, pp. 888--891.
        \href{https://doi.org/10.1145/3238147.3240481}{doi:10.1145/3238147.3240481}
  \item B. Farias, R. Menezes, E. B. de Lima Filho, Y. Sun, L. C. Cordeiro.
        ``ESBMC-Python: A Bounded Model Checker for Python Programs''.
        \emph{ISSTA 2024}, pp. 1836--1840.
        \href{https://doi.org/10.1145/3650212.3685304}{doi:10.1145/3650212.3685304}
  \item N. Tihanyi, Y. Charalambous, R. Jain, M. A. Ferrag, L. C. Cordeiro.
        ``A New Era in Software Security: Towards Self-Healing Software via
        Large Language Models and Formal Verification''. \emph{AST@ICSE 2025},
        pp. 136--147.
        \href{https://doi.org/10.1109/AST66626.2025.00020}{doi:10.1109/AST66626.2025.00020}
  \item C. Baier, J.-P. Katoen. \emph{Principles of Model Checking}. MIT Press,
        2008. ISBN 978-0-262-02649-9.
  \item A. Burns, A. J. Wellings. \emph{Real-Time Systems and Programming
        Languages: Ada, Real-Time Java and C/Real-Time POSIX}, 4th edition.
        Addison-Wesley, 2009. ISBN 978-0-321-41745-9. Source of the
        chemical-process case study in Exercise~\ref{ex:plant}.
\end{enumerate}

\end{document}
